% Overleaf-ready paper skeleton and front matter.
% Compile with pdfLaTeX. No bibliography entries are included because all
% literature references must be independently verified before citation.
%
% Set \planningappendixfalse for a clean conference-submission skeleton.
\documentclass[10pt,twocolumn]{article}

\usepackage[T1]{fontenc}
\usepackage[utf8]{inputenc}
\usepackage{lmodern}
\usepackage[a4paper,margin=18mm,columnsep=6mm]{geometry}
\usepackage{microtype}
\usepackage{booktabs}
\usepackage{tabularx}
\usepackage{array}
\usepackage{enumitem}
\usepackage{xcolor}
\usepackage{textcomp}
\usepackage{amsmath}
\usepackage{hyperref}
\usepackage{caption}

\hypersetup{
  colorlinks=true,
  linkcolor=blue!55!black,
  urlcolor=blue!55!black,
  citecolor=blue!55!black,
  pdftitle={When Better Scenario Dependence Does Not Pay},
  pdfauthor={Author placeholder}
}

\setlength{\parindent}{0.9em}
\setlength{\parskip}{0pt}
\setlist{nosep,leftmargin=*}
\captionsetup{font=small,labelfont=bf}

\newif\ifplanningappendix
\planningappendixtrue % Change to \planningappendixfalse for submission.

\newcommand{\status}[1]{\texttt{\detokenize{#1}}}
\newcommand{\tverify}{\textsc{To verify}}
\newcommand{\claimbox}[1]{%
  \begin{center}
  \fbox{\begin{minipage}{0.94\columnwidth}\small #1\end{minipage}}
  \end{center}}

\title{\textbf{When Better Scenario Dependence Does Not Pay:}\\
A Feasibility-Restricted Evaluation of Coherent Forecast Uncertainty for
Industrial Flexibility Scheduling}

\author{
  Author Name(s)\\
  \small Affiliation / Department / Institution\\
  \small City, Country\\
  \small \texttt{corresponding.author@example.org}
}
\date{}

\begin{document}
\maketitle

\begin{center}
\small\textbf{Intended venue type:} German/European applied energy informatics,
industrial flexibility, or power-systems conference (e.g., The smarter E
Europe, pro flex Dresden, Energy Informatics, or an ETG/VDE energy event).
\end{center}

\begin{abstract}
Industrial flexibility scheduling increasingly uses probabilistic forecasts,
yet improved statistical representation of temporal dependence need not yield
better operational decisions. Existing evaluations often separate scenario
quality from downstream feasibility and economic value, leaving the practical
benefit of coherent daily trajectories unclear. We conduct a leakage-safe,
feasibility-restricted counterfactual study of day-ahead scheduling with a
frozen virtual production-buffer asset. Deterministic LightGBM scheduling is
compared with risk-neutral schedules driven by independent interval residuals
and coherent whole-day residual trajectories. The empirical decision analysis
uses 273 complete 15-minute KIT SPARK days from 6 January to 1 December 2024;
191 days are feasible for all three policies. Supporting forecast and scenario
diagnostics draw on SPARK positive facility demand and Honda signed net grid
exchange, whose targets are treated as physically distinct. On common-feasible
days, deterministic scheduling is cheaper than coherent-scenario scheduling
by 0.191191 euro/day (95\% day-bootstrap interval: 0.131960--0.271593 euro/day
in magnitude). The coherent-versus-independent difference is statistically
unresolved: independent minus coherent equals -0.047240 euro/day (95\%
interval: -0.122619 to 0.009287). Nevertheless, coherent scenarios reduce
modelled zero-export-cap violations from 82 deterministic days and 16
independent-scenario days to 11 days. Thus, better temporal dependence can
improve modelled feasibility without improving mean synthetic counterfactual
wholesale cost. Results are exploratory, depend on a hypothetical
$\pm5$ kW/$\pm10$ kWh flexibility asset and an assumed post-publication
decision time, and constitute neither plant validation nor customer savings.
\end{abstract}

\noindent\textbf{Keywords:} industrial flexibility; probabilistic forecasting;
scenario coherence; decision value; day-ahead scheduling; forecast uncertainty;
feasibility; signed grid exchange

\claimbox{\textbf{One-sentence contribution.} This paper presents a
leakage-safe, feasibility-restricted counterfactual study showing that coherent
whole-day residual scenarios can reduce constraint violations without
improving mean day-ahead decision cost, demonstrating that temporal scenario
dependence is not equivalent to decision value.}

\section{Introduction}
\label{sec:introduction}

Industrial energy-management studies commonly evaluate point or probabilistic
forecasts using statistical scores, while the downstream scheduling problem is
evaluated separately or under weak physical assumptions. This separation can
obscure whether an improvement in the temporal structure of forecast scenarios
changes an implementable decision. The question addressed here is therefore
narrow and falsifiable: \emph{Does better temporal coherence in probabilistic
load or net-exchange scenarios improve day-ahead industrial flexibility
decisions?}

The analysis deliberately preserves a negative result. Whole-day residual
trajectories improve dependence diagnostics and substantially reduce modelled
zero-export-cap violations, but they do not lower mean synthetic
counterfactual wholesale cost under the frozen decision problem. The paper is
not framed as a forecasting-model comparison, a battery study, or validation
of plant operations.

\subsection{Contributions}

\begin{itemize}
  \item A leakage-safe, feasibility-restricted decision-value evaluation
  across positive demand and signed net-exchange targets.
  \item A reproducible comparison of deterministic, independent-scenario, and
  coherent-scenario scheduling under a frozen counterfactual flexibility asset.
  \item A bounded negative result showing that improved temporal scenario
  dependence reduces constraint violations but does not improve mean decision
  cost.
\end{itemize}

\claimbox{\textbf{Claims and non-claims.}

\textbf{Allowed claims:} coherent scenarios reduce modelled violation
frequency; deterministic scheduling remains cheaper in mean synthetic
counterfactual wholesale cost; coherent-versus-independent feasible-day cost
is inconclusive; the explicit feasibility correction does not reverse the cost
ranking; and the result is bounded to SPARK, the frozen F1 asset, and the
assumed timing convention.

\textbf{Not claimed:} actual customer savings, trading profit, plant
validation, Honda economic value, general industrial-flexibility value,
battery value, trading performance, or universal superiority or inferiority of
coherent scenarios.}

\section{Related Work and Research Gap}
\label{sec:related}

The final paper should organize verified literature into six strands:
(i) industrial electricity and net-exchange forecasting; (ii) probabilistic
forecasting and quantile calibration; (iii) scenario generation and temporal
dependence; (iv) decision-focused learning and predict-then-optimize methods;
(v) industrial demand-side flexibility; and (vi) electricity-market timing,
settlement, and information availability. The synthesis should distinguish
proper marginal calibration from the validity of joint daily trajectories.

The intended gap statement is modest: forecast quality is frequently assessed
separately from downstream decision value, and relatively little evidence
shows when improved temporal dependence changes a physically constrained,
day-ahead industrial schedule. All literature claims and citations remain
\tverify{} until checked against primary sources; Section~\ref{app:references}
defines the reference plan without inventing citations.

\section{Data and Provenance}
\label{sec:data}

\subsection{Dataset semantics}

The empirical foundation comprises two industrial energy datasets with
different physical targets. KIT SPARK records positive facility electrical
demand. The modelling target is a 15-minute arithmetic mean constructed from
raw measurements sampled every five seconds and spans 2021--2024. Honda R\&D
Europe in Offenbach records 15-minute signed exchange with the public grid for
2018--2023, with primary forecasting labels concentrated in 2020--2023.
Positive Honda values denote import and negative values export. The Honda
target already nets onsite PV and CHP and is therefore not gross demand.
Honda gross-demand reconstruction has not been independently validated and is
not an approved forecasting target.

The physical distinction is consequential. SPARK describes positive facility
consumption, whereas Honda describes bidirectional net exchange after onsite
generation. Raw kW forecast errors across the two sites therefore do not measure
the same quantity and must not be compared as if they did. Honda is used to
evaluate scenario construction and forecasting generality, but the economic
decision-value replication in this paper is restricted to SPARK.

\subsection{Measurement quality and target provenance}

SPARK contains substantial raw missingness. Table~\ref{tab:spark-quality}
reports valid five-second samples, availability, and the longest continuous gap
by year. Availability falls from 78.7746\% in 2021 to 68.3132\% in 2024, while
the longest gap expands from less than two hours to almost 30 days. These values
motivate explicit interval-level eligibility rather than treating the archive
as a complete regular series.

\begin{table*}[t]
\centering
\caption{SPARK raw-measurement availability and longest missing gap.}
\label{tab:spark-quality}
\small
\begin{tabular}{lrrl}
\toprule
Year & Valid source samples & Availability & Longest missing gap\\
\midrule
2021 & 4,968,471 & 78.7746\% & 1 h 34 min 40 s\\
2022 & 5,023,226 & 79.6427\% & 24 h 53 min 30 s\\
2023 & 4,520,517 & 71.6723\% & 13 d 16 h 49 min\\
2024 & 4,320,457 & 68.3132\% & 29 d 16 h 16 min\\
\bottomrule
\end{tabular}
\end{table*}

SPARK target construction uses an empirically supported 75\% raw-coverage
threshold within each 15-minute interval. Table~\ref{tab:aggregation-error}
shows that aggregation error rises markedly around 70\% coverage: mean absolute
error increases from 0.0060 to 0.0416 kW, and the P99 error from 0.0495 to
0.2687 kW. The 75\% rule is consequently a site- and archive-specific quality
decision, not a universal threshold for industrial data.

\begin{table*}[t]
\centering
\caption{Empirical SPARK aggregation error by within-interval raw coverage.}
\label{tab:aggregation-error}
\small
\begin{tabular}{lrrr}
\toprule
Coverage & Mean absolute aggregation error & P95 absolute error & P99 absolute error\\
 & (kW) & (kW) & (kW)\\
\midrule
75\% & 0.0060 & 0.0212 & 0.0495\\
Approximately 70\% & 0.0416 & 0.1364 & 0.2687\\
\bottomrule
\end{tabular}
\end{table*}

Honda target construction is provenance-aware. Primary labels retain only
original observations; author-imputed or filled values are excluded. Corrected
non-imputed values belong only to a separately labelled sensitivity policy.
Unknown provenance and unresolved missing intervals are also excluded from
primary labels. This policy prevents filled values from being presented as
independent measurements and preserves the signed import/export target.

\section{Forecasting Models and Evaluation Protocol}
\label{sec:forecasting}

Forecasts are issued at 12:00 Europe/Berlin on D-1 for SPARK and 13:00 for
Honda. Each forecast covers every physical 15-minute interval of local day D.
The use of UTC as the unique physical key preserves 92 intervals on the spring
DST transition, 96 on an ordinary day, and 100 on the autumn transition. Only
measurements whose physical intervals have completed by issuance are available
to a model.

All model development is chronological. The SPARK deterministic benchmark
uses 2021--2022 for training, 2023 for validation, and 2024 for testing. Honda
uses rolling development folds with 2020 training/2021 validation and
2020--2021 training/2022 validation, followed by refitting on 2020--2022 and a
single 2023 evaluation. No future measured load, grid exchange, PV, CHP,
weather, or target-derived regime variable is admitted. These timing rules are
also applied to historical analogues and feature construction.

For SPARK, deterministic LightGBM is the principal scheduling forecast. B3 is
a robust recent-week analogue and B5 a selected historical-baseline ensemble.
Table~\ref{tab:spark-forecast} reports the frozen 2024 comparison. LightGBM has
the lowest MAE and RMSE and full prediction coverage; B5 has 96.19\% coverage.
These frozen forecasts enter the decision study without further model selection
or tuning.

\begin{table}[t]
\centering
\caption{SPARK 2024 deterministic forecasting results.}
\label{tab:spark-forecast}
\small
\begin{tabular}{lrrr}
\toprule
Method & MAE (kW) & RMSE (kW) & Coverage\\
\midrule
B3 & 2.7096 & 4.2319 & 100\%\\
B5 & 2.5713 & 4.1518 & 96.19\%\\
LightGBM & 2.3520 & 3.4659 & 100\%\\
\bottomrule
\end{tabular}
\end{table}

The SPARK probabilistic model estimates q05, q10, q25, q50, q75, q90, and q95
with quantile LightGBM. A development-selected global calibration produces 2024
WIS 1.4743 and empirical 50\%, 80\%, and 90\% coverages of 50.06\%, 80.44\%,
and 89.91\%. Holiday calibration remains weak: WIS is 2.1746, while 80\% and
90\% coverage fall to 50.35\% and 68.29\%. These marginal quantile results do
not, by themselves, establish a coherent joint daily distribution.

Honda provides a bounded cross-dataset forecasting benchmark. As shown in
Table~\ref{tab:honda-forecast}, LightGBM M1-CHP improves aggregate MAE over B5
by 4.967 kW but does not improve export recall. Export recall is 29.752\% for
LightGBM and 30.809\% for B5. Zero-shot TimesFM 3.0 and Chronos-2 are not
competitive in this implementation. This result concerns two particular
checkpoints in zero-shot mode; it is not evidence that foundation models are
universally inferior or that adaptation cannot help.

\begin{table*}[t]
\centering
\caption{Honda 2023 signed-grid-exchange forecasting results. Positive targets
denote import and negative targets export.}
\label{tab:honda-forecast}
\small
\begin{tabular}{lrrrr}
\toprule
Model & MAE (kW) & RMSE (kW) & Bias (kW) & Export recall\\
\midrule
B5 & 76.480 & 101.338 & +6.649 & 30.809\%\\
LightGBM M1-CHP & 71.513 & 96.974 & +9.408 & 29.752\%\\
TimesFM 3.0 zero-shot & 95.431 & 129.221 & -0.040 & 11.219\%\\
Chronos-2 zero-shot & 103.672 & 139.959 & -6.184 & 11.543\%\\
\bottomrule
\end{tabular}
\end{table*}

\section{Scenario Generation}
\label{sec:scenarios}

For observation $y_{d,t}$ and frozen point forecast $\widehat{y}_{d,t}$, the
residual is
\begin{equation}
  e_{d,t}=y_{d,t}-\widehat{y}_{d,t}.
\end{equation}
Residual donors are restricted to past or out-of-sample development residuals;
test residuals are never sampled to construct a test-day scenario. S1 generates
each interval by independently sampling residual information, thereby
preserving a marginal residual distribution while disrupting within-day serial
dependence. S2 samples a whole-day residual trajectory and adds it to the point
forecast, preserving the donor day's temporal ordering. S3 conditions donor
selection on pre-specified information, but was not selected because it did not
pass a robust across-dataset selection rule.

Tables~\ref{tab:honda-scenarios} and~\ref{tab:spark-scenarios} compare S1 and S2
at 50 scenarios. Lower scores are preferable. S2 produces modestly lower energy
scores but much lower variogram scores and ACF RMSE on both datasets. The
strongest improvement is therefore in dependence representation rather than in
the overall multivariate distance measured by the energy score.

\begin{table}[t]
\centering
\caption{Honda 2022 scenario diagnostics at 50 scenarios.}
\label{tab:honda-scenarios}
\small
\begin{tabular}{lrr}
\toprule
Metric & S1 & S2\\
\midrule
Energy score & 64.352 & 63.317\\
Variogram score & 25.040 & 8.256\\
ACF RMSE & 0.498 & 0.211\\
\bottomrule
\end{tabular}
\end{table}

\begin{table}[t]
\centering
\caption{SPARK rolling-2023 scenario diagnostics.}
\label{tab:spark-scenarios}
\small
\begin{tabular}{lrr}
\toprule
Metric & S1 & S2\\
\midrule
Energy score & 2.394 & 2.371\\
Variogram score & 1.238 & 0.402\\
ACF RMSE & 0.538 & 0.185\\
\bottomrule
\end{tabular}
\end{table}

S2 is not claimed to be fully marginally calibrated. Table~\ref{tab:coverage}
shows undercoverage at both nominal levels: 80\% intervals cover 75.21\% for
Honda and 69.74\% for SPARK, while 90\% intervals cover 85.56\% and 81.29\%.
Consequently, S2 improves dependence metrics but does not solve marginal
calibration. Better dependence is an input property whose decision value must
be tested, not assumed.

\begin{table}[t]
\centering
\caption{Marginal coverage limitation of S2 scenarios.}
\label{tab:coverage}
\small
\begin{tabular}{lrr}
\toprule
Dataset & S2 80\% coverage & S2 90\% coverage\\
\midrule
Honda & 75.21\% & 85.56\%\\
SPARK & 69.74\% & 81.29\%\\
\bottomrule
\end{tabular}
\end{table}

\section{Counterfactual Flexibility and Market-Timing Framework}
\label{sec:flexibility}

\subsection{Frozen F1 work-buffer asset}

The decision model uses a synthetic counterfactual electrical-work buffer F1.
For baseline exchange or demand $b_t$, committed control $u_t$, and interval
duration $\Delta t=0.25$ h, the post-control grid trajectory $G_t$ and work
state $x_t$ are
\begin{align}
  G_t &= b_t+u_t,\\
  x_{t+1} &= x_t+u_t\Delta t.
\end{align}
The frozen physical bounds are
\begin{align}
  -P_{\mathrm{down}} &\leq u_t \leq P_{\mathrm{up}},\\
  -E_{\mathrm{backlog}} &\leq x_t \leq E_{\mathrm{advance}},
\end{align}
with $P_{\mathrm{up}}=P_{\mathrm{down}}=5$ kW and
$E_{\mathrm{advance}}=E_{\mathrm{backlog}}=10$ kWh. Exact daily energy
conservation follows from terminal recovery $x_T=x_0$. The model imposes a
zero-export cap, no import cap, and one shared first-stage non-anticipative
schedule across all scenarios.

F1 is neither an identified machine nor a measured production line. It is not a
battery or a thermal system, and its capacities are not measured equipment
parameters. Thermal-buffer F2 and BESS formulations are outside the primary
decision-value experiment. For Honda, any future use of F1 would be an additive
hypothetical load overlay on signed exchange; PV and CHP must not be subtracted
again. No Honda economic result is reported here.

\subsection{Market information and settlement}

SPARK's load forecast is issued at 12:00 Europe/Berlin on D-1. Scheduling is
assumed at 12:55, after the relevant day-ahead price vector has become known.
The historical price file does not contain observation-level publication
timestamps, so the information pattern is labelled
\status{POST_PUBLICATION_TIMING_ASSUMED}. The original experiment remains
\status{EX_POST_PRICE_AWARE}; neither label establishes real-time operational
validation. A later decision time does not permit the 12:00 load forecast to be
refreshed with post-issuance measurements.

The implementable schedule is sealed before actual demand is accessed. If the
realized post-control trajectory $g_t$ violates the zero-export cap, the
secondary all-day estimand applies
\begin{align}
  c_t &= \max(0,-g_t),\\
  g^{\mathrm{corrected}}_t &= g_t+c_t.
\end{align}
This is the unique componentwise minimum nonnegative intervention needed to
satisfy $g^{\mathrm{corrected}}_t\geq0$. It equals, and therefore cannot exceed,
the violation amount. Because the intervention is zero whenever $g_t\geq0$ and
only removes negative exchange, it cannot create additional negative-price
consumption arbitrage. The committed F1 schedule is never changed, clipped, or
reoptimized, and the correction lies outside the F1 work-buffer state. It is an
explicit settlement sensitivity rather than hidden recourse.

\section{Feasibility-Restricted Decision-Value Evaluation}
\label{sec:evaluation}

The evaluation population contains 273 complete 96-interval SPARK days from
6 January through 1 December 2024. Spring and autumn DST days are valid in the
general protocol, but this replication population contains ordinary 96-interval
days only. A policy-day is feasible when every realized interval satisfies the
zero-export cap within $10^{-7}$ kW. A common-feasible day is one for which P2,
P3, and P4 are all feasible; 191 days satisfy this intersection.

P2 schedules against deterministic LightGBM, P3 against S1 independent
scenarios, and P4 against S2 coherent scenarios. The same price vector, F1
asset, terminal condition, solver convention, and realized settlement are used
for all policies. Exactly three comparisons were pre-registered: P2$-$P4,
P3$-$P4, and P2$-$P3. For cost difference A$-$B, a positive result means that B
is cheaper and a negative result that A is cheaper.

The \status{PRE_REGISTERED_PRIMARY} estimand uses only the 191 common-feasible
days. Uncertainty is quantified with 2,000 paired whole-day bootstrap resamples
using seed 20261029. Seven-day circular moving-block intervals are a temporal
dependence sensitivity, and Holm adjustment covers the three primary tests.
The corrected 273-day comparisons are
\status{PRE_SPECIFIED_SECONDARY}. Nominal uncorrected all-day comparisons,
violation diagnostics, and binding-constraint diagnostics are
\status{EXPLORATORY_DIAGNOSTIC}. These labels prevent the all-day sensitivity
from replacing the primary common-feasible estimand after outcomes are known.

\section{Results}
\label{sec:results}

\subsection{Policy-level outcomes}

Table~\ref{tab:policy} reports the policy-level results. The values called cost
reductions in the caption are synthetic counterfactual wholesale-cost
differences relative to no flexibility, not customer savings or trading profit.
P2 has the greatest nominal, common-feasible, and corrected reduction in cost.
P4 has the fewest modelled violation days but the smallest mean cost reduction,
so improved feasibility does not coincide with improved mean economics.

\begin{table*}[t]
\centering
\caption{Policy-level results. All monetary values are synthetic counterfactual
wholesale-cost reductions, not actual savings.}
\label{tab:policy}
\small
\begin{tabular}{lrrrr}
\toprule
Policy & Nominal cost reduction & Violation days & Common-feasible cost reduction & Corrected cost reduction\\
 & (euro/day) & & (euro/day) & (euro/day)\\
\midrule
P2 deterministic LightGBM & 1.999617 & 82 & 2.326373 & 1.982463\\
P3 independent scenarios & 1.800671 & 16 & 2.182422 & 1.799640\\
P4 coherent scenarios & 1.747536 & 11 & 2.135182 & 1.746470\\
\bottomrule
\end{tabular}
\end{table*}

\subsection{Primary common-feasible comparisons}

Table~\ref{tab:primary-comparisons} contains the three
\status{PRE_REGISTERED_PRIMARY} comparisons. P2 is robustly cheaper than P4:
the mean P2$-$P4 difference is -0.191191 euro/day, and both day and block
intervals remain below zero. P2 is likewise cheaper than P3 by 0.143951
euro/day in magnitude. By contrast, P3$-$P4 equals -0.047240 euro/day and both
uncertainty intervals include zero; its Holm-adjusted $p$-value is 0.131.
Although P4 is cheaper than P3 on 65.97\% of days, larger adverse P4 outcomes
raise its mean cost. The coherent-versus-independent mean comparison is
therefore inconclusive, not evidence of average superiority in either direction.

\begin{table*}[t]
\centering
\caption{Primary common-feasible paired comparisons. Positive A$-$B means B is
cheaper; negative values mean A is cheaper.}
\label{tab:primary-comparisons}
\scriptsize
\begin{tabular}{lrrrrrr}
\toprule
Comparison & Days & Mean A$-$B & 95\% day-bootstrap CI & 7-day block CI & B daily win fraction & Holm-adjusted $p$\\
 & & (euro/day) & (euro/day) & (euro/day) & & \\
\midrule
P2$-$P4 & 191 & -0.191191 & [-0.271593, -0.131960] & [-0.313554, -0.113278] & 0.0052 & 0.000\\
P3$-$P4 & 191 & -0.047240 & [-0.122619, +0.009287] & [-0.135098, +0.005784] & 0.6597 & 0.131\\
P2$-$P3 & 191 & -0.143951 & [-0.169757, -0.120702] & [-0.188752, -0.104662] & 0.0000 & 0.000\\
\bottomrule
\end{tabular}
\end{table*}

\subsection{Corrected all-day sensitivity}

The \status{PRE_SPECIFIED_SECONDARY} corrected comparison in
Table~\ref{tab:corrected-comparisons} uses all 273 days. The correction does not
reverse the policy ranking. P2 remains cheaper than both scenario policies. P3
is cheaper than P4 by 0.053170 euro/day in mean corrected cost, with both day
and block intervals below zero. This secondary result is conditional on the
external correction and does not supersede the primary common-feasible finding.

\begin{table*}[t]
\centering
\caption{Corrected all-day paired comparisons
(\status{PRE_SPECIFIED_SECONDARY}).}
\label{tab:corrected-comparisons}
\scriptsize
\begin{tabular}{lrrrrr}
\toprule
Comparison & Days & Mean A$-$B & 95\% day-bootstrap CI & 7-day block CI & B daily win fraction\\
 & & (euro/day) & (euro/day) & (euro/day) & \\
\midrule
P2$-$P4 & 273 & -0.235992 & [-0.300416, -0.184113] & [-0.321147, -0.172512] & 0.0330\\
P3$-$P4 & 273 & -0.053170 & [-0.109739, -0.011935] & [-0.113356, -0.011379] & 0.6154\\
P2$-$P3 & 273 & -0.182822 & [-0.209122, -0.157245] & [-0.221738, -0.145244] & 0.0293\\
\bottomrule
\end{tabular}
\end{table*}

\subsection{Feasibility and overlap diagnostics}

Table~\ref{tab:feasibility} quantifies modelled zero-export-cap violations.
P4 reduces violation days by 71, or 86.6\%, relative to P2 and by five, or
31.25\%, relative to P3. Its maximum violation power is also the smallest.
However, among its 11 violating days, mean violation energy exceeds P3's mean
over 16 violating days. The diagnostic supports a reduction in violation
frequency, not a general claim of lower violation severity or economic value.

\begin{table*}[t]
\centering
\caption{Feasibility diagnostics (\status{EXPLORATORY_DIAGNOSTIC}).}
\label{tab:feasibility}
\scriptsize
\begin{tabular}{lrrrrrrr}
\toprule
Policy & Compliant days & Violation days & Violation fraction & Mean energy & P90 energy & Maximum energy & Maximum power\\
 & & & & (kWh) & (kWh) & (kWh) & (kW)\\
\midrule
P2 & 191 & 82 & 30.04\% & 0.7093 & 2.0154 & 9.4713 & 1.1259\\
P3 & 257 & 16 & 5.86\% & 0.2174 & 0.3281 & 2.0182 & 0.5017\\
P4 & 262 & 11 & 4.03\% & 0.3495 & 0.5171 & 2.0295 & 0.4599\\
\bottomrule
\end{tabular}
\end{table*}

Violation sets are strongly nested (Table~\ref{tab:overlap}). All 11 P4
violation days also violate under P2. Sixty-five days violate only under P2,
whereas no P3- or P4-only violation occurs without P2. The overlap explains why
the common-feasible intersection equals P2's 191 compliant days and confirms
that P4's reduced violation count is not produced by a disjoint selection of
days.

\begin{table}[t]
\centering
\caption{Overlap of modelled violation days
(\status{EXPLORATORY_DIAGNOSTIC}).}
\label{tab:overlap}
\small
\begin{tabular}{lr}
\toprule
Violation pattern & Days\\
\midrule
No policy violates & 191\\
P2 only & 65\\
P2 and P3 & 6\\
P2 and P4 & 1\\
All three & 10\\
P3 or P4 without P2 & 0\\
\bottomrule
\end{tabular}
\end{table}

Finally, Table~\ref{tab:binding} shows that F1 power and energy-state limits
bind on most days under every policy. Power limits bind on all P2 days and all
but two P3 days; energy-state limits bind on at least 251 days. Binding is thus
too common to identify a useful subgroup in which coherence pays. The result
does not establish that the P4 advantage is concentrated on binding days.

\begin{table}[t]
\centering
\caption{Binding-constraint diagnostics
(\status{EXPLORATORY_DIAGNOSTIC}).}
\label{tab:binding}
\small
\begin{tabular}{lrrr}
\toprule
Policy & Power-limit & Energy-state & Grid-cap\\
 & binding days & binding days & violation days\\
\midrule
P2 & 273 & 273 & 82\\
P3 & 271 & 253 & 16\\
P4 & 266 & 251 & 11\\
\bottomrule
\end{tabular}
\end{table}

\section{Discussion}
\label{sec:discussion}

\subsection{Mechanism analysis: why coherence did not pay}
\label{sec:mechanism}

The empirical result is consistent with the mathematical structure of the
decision problem. Under a shared first-stage schedule and a risk-neutral
expected energy-cost objective, much of the expected cost is driven by marginal
expectations. Temporal dependence can alter the optimum when uncertainty enters
through time-coupled state constraints, joint import/export caps, daily peak
terms, nonlinear tariffs, CVaR, chance constraints, or explicitly timed
recourse. It need not materially alter an expected-cost schedule merely because
its scenario trajectories have more realistic autocorrelation.

The F1 state equation creates temporal coupling, and the zero-export cap makes
the joint trajectory relevant to feasibility. S2 accordingly reduces the
frequency of modelled cap violations. Yet power and state bounds bind on nearly
all days, limiting the schedules available to each policy. P4's response to
joint residual persistence is more conservative: it protects against coherent
adverse trajectories but can shift flexible work into less favourable price
intervals. In this experiment the reduction in infeasibility does not offset
that mean cost effect.

The distinction is visible in the paired statistics. P4 is cheaper than P3 on
65.97\% of common-feasible days, but a smaller number of adverse outcomes makes
its average cost higher. Because the P3$-$P4 day-bootstrap and 7-day block
intervals both cross zero and the Holm-adjusted $p$-value is 0.131, the evidence
does not establish that S2 is better or worse on average. It establishes a
trade-off: better temporal dependence and fewer modelled violations coexist
with no demonstrated improvement in mean synthetic counterfactual wholesale
cost.

\emph{Improved temporal dependence is not equivalent to improved decision
value.} Energy score, variogram score, and autocorrelation diagnostics remain
necessary checks on scenario quality, but none establishes that a particular
uncertainty representation is aligned with an economic objective. Decision
readiness additionally requires a fixed information pattern, physical
feasibility, identical settlement, and paired downstream evaluation.

This mechanism is bounded to the frozen SPARK experiment. Different conclusions
could arise with asymmetric import/export remuneration, nonseparable peak
charges, scenario-wise hard constraints, a risk-aware objective, a different
physical asset, or validated recourse. Those possibilities are hypotheses for
separate pre-specified studies, not post-hoc explanations that alter the present
negative result.

The finding is compatible with the broader possibility that probabilistic
scheduling produces more conservative schedules and, in some stochastic
settings, lower realized profit than deterministic scheduling. That compatibility
is an interpretation to be assessed against verified literature, not a new
empirical result of this paper. Here, the observed trade-off is specific:
coherence reduces modelled violations, but its conservative schedule does not
lower average cost under the frozen counterfactual.

\section{Practical Implications}
\label{sec:implications}

\textbf{Industrial energy managers.} Coherent scenario scheduling should not
be adopted solely because it improves probabilistic forecast diagnostics. The
relevant question is whether the actual flexibility resource has state coupling,
binding caps, recovery constraints, nonlinear costs, or risk measures that make
joint trajectories economically consequential. A strong deterministic schedule
should remain a mandatory baseline, and feasibility must be assessed under
realized operation rather than forecast trajectories alone.

\textbf{Energy-analytics teams.} Model evaluation should connect MAE, RMSE,
WIS, energy score, and variogram score to the actual downstream objective.
Marginal calibration and temporal coherence should be reported separately.
Cost tables should be accompanied by prediction coverage, common-day support,
solver feasibility, and constraint-violation diagnostics so that a favorable
mean is not produced by different or easier evaluation populations.

\textbf{Researchers.} When the scientific objective is scheduling or market
participation, decision value should be evaluated directly. Policy comparisons
should be pre-registered, preserve identical physical and market assumptions,
and use a common-feasible-day estimand when infeasibility differs. Negative
results should be reported transparently: they can identify which uncertainty
properties the optimization does, and does not, reward.

\section{Limitations}
\label{sec:limitations}

The following limitations define the scope of the evidence.

\begin{enumerate}
  \item The F1 resource is hypothetical, not an identified industrial process.
  \item Historical price publication timestamps are absent; 12:55 is assumed
  from documentary convention.
  \item SPARK 2024 was previously exposed, so the replication is not independent
  confirmation.
  \item The correction is an ex-post external feasibility ledger, not an
  implementable plant action.
  \item Common-feasible filtering changes the evaluated population and may
  select less difficult days.
  \item Only SPARK is economically evaluated; no Honda economic inference is
  made.
  \item Every evaluated day has 96 intervals; the experiment provides no
  DST-transition economic result.
  \item Wholesale exposure is not a customer tariff or total electricity bill.
  \item The scenario generator improves dependence but is not fully marginally
  calibrated.
  \item Two public sites cannot establish general industrial external validity.
  \item The zero-shot foundation-model result is limited to Chronos-2 and
  TimesFM 3.0 and does not evaluate fine-tuning or domain adaptation.
\end{enumerate}

\section{Future Work}
\label{sec:future}

Further work should be restricted to extensions that directly address an
identified limitation or decision mechanism:

\begin{enumerate}
  \item add a third public industrial dataset with a different flexibility
  structure;
  \item add a transparent, publicly documented German industrial tariff or
  network-charge scenario;
  \item develop a feasibility-aware or decision-aware scenario generator and
  compare it with S1 and S2;
  \item decompose the deterministic-versus-coherent cost difference into
  conservatism cost, avoided-violation value, and timing effects;
  \item test coherent scenarios where state coupling is stronger, including
  multi-day storage, demand charges, or explicitly timed recourse decisions;
  \item obtain authoritative Honda price and publication-timestamp data before
  attempting any Honda economic evaluation.
\end{enumerate}

These priorities exclude additional zero-shot checkpoint surveys, further
regime-calibration variants, unconstrained capacity sweeps, battery modelling
without a real asset, and Honda economics without verified prices and timing.

\section{Conclusion}
\label{sec:conclusion}

This study asked whether a better representation of temporal dependence in
probabilistic scenarios improves day-ahead industrial flexibility decisions.
It compared a frozen deterministic LightGBM schedule with risk-neutral schedules
based on independent interval residuals and coherent whole-day residual
trajectories. All policies used the same hypothetical $\pm5$ kW/$\pm10$ kWh
electrical-work buffer, wholesale-price vector, terminal recovery rule, and
zero-export cap. The primary comparison was restricted to 191 SPARK days that
were feasible under every policy.

The main economic result is negative. Deterministic scheduling was cheaper than
coherent-scenario scheduling by 0.191191 euro/day, with day-bootstrap and 7-day
block intervals wholly below zero. It was also cheaper than independent-scenario
scheduling. The coherent-versus-independent difference was statistically
unresolved. In contrast, coherence had a clear feasibility effect: modelled
zero-export-cap violations fell from 82 deterministic days and 16 independent-
scenario days to 11 coherent-scenario days. An explicit all-day correction did
not reverse the cost ranking.

The practical implication is that improved dependence diagnostics cannot
substitute for downstream evaluation. Coherent scenarios can make a schedule
more conservative and reduce modelled violations without lowering its average
cost. This conclusion is bounded to one public dataset, a hypothetical asset,
wholesale exposure, imperfectly calibrated scenarios, and
\status{POST_PUBLICATION_TIMING_ASSUMED}; it is neither plant validation nor
evidence of customer savings or trading profit. Probabilistic scenario
sophistication should be adopted only when its uncertainty representation
matches the economic structure of the downstream decision.

\section{Reproducibility and Data Availability}
\label{sec:reproducibility}

The replication configuration was frozen before outcome evaluation. Two
complete runs produced byte-identical deterministic outputs, and the full test
suite passed 233 of 233 tests, including 20 tests introduced specifically for
the replication. Original Stage 1--3 artifacts were unchanged. No forecasts
were retrained, no scenarios were regenerated, and the read-only SPARK
repository remained clean. Table~\ref{tab:artifact-hashes} records the principal
artifact fingerprints.

\begin{table*}[t]
\centering
\caption{Reproducibility artifacts and SHA-256 fingerprints.}
\label{tab:artifact-hashes}
\scriptsize
\begin{tabularx}{\textwidth}{>{\ttfamily}X >{\ttfamily}p{0.46\textwidth}}
\toprule
\normalfont Artifact & \normalfont SHA-256\\
\midrule
config/decision\_value\_replication/stage3\_feasibility\_replication\_frozen.json & f00a24c55072cce5ec9399d506231d6f7fe67d29ae50448a8ea9ba1d3e39c0c3\\
src/honda/decision\_value\_replication/replication.py & 2dfc078f061ad9a282557f7af653f79b19f697f4f230d2da5ead9d61927cd4b0\\
results/decision\_value\_replication/replication\_manifest.json & 88bc917eed7cce7601ba2837d8884fca798e33a280e1f7702f1b6dc3ff1b7949\\
results/decision\_value\_replication/reproducibility\_check.json & 2015c726d850d73ec8cd1a520a723dfa7f69b4a2a91040dfa02b65edb9aa14e5\\
results/decision\_value\_replication/replication\_policy\_summary.csv & 7e4c5d5d6d8f97913a69d697033d5ad2a7266b348684adc3c9e1b9127bcc49fa\\
results/decision\_value\_replication/primary\_common\_feasible\_comparisons.csv & 30819f4ecb7ee452bbdec74776504818ac20f59f559d723d07057127a129083a\\
results/decision\_value\_replication/secondary\_corrected\_comparisons.csv & 577a7ceea27781051f2df78c9985a04acf245fb8b0a535abaa84582de48f4160\\
results/decision\_value\_replication/feasibility\_diagnostics.csv & ccf4291e36de30a315db066342603ae5c7234160ffcc254f6230710a1256c995\\
results/decision\_value\_replication/binding\_constraint\_diagnostics.csv & 6278be55e527c15fa257e6e1efd9c15f318d66b916d74af6c4761ca1c0fb222e\\
\bottomrule
\end{tabularx}
\end{table*}

The study uses the public KIT SPARK and Honda R\&D Europe, Offenbach datasets.
Code, configuration, and derived result artifacts are intended for release
subject to repository, dataset, and model-licensing requirements. This statement
does not promise redistribution of source data or artifacts whose licenses have
not been independently confirmed.

\section{Final Claims and Non-Claims}
\label{sec:claims}

\claimbox{\textbf{Allowed claims}
\begin{itemize}
  \item Coherent scenarios reduce modelled zero-export-cap violation frequency.
  \item Deterministic LightGBM remains cheaper in mean synthetic counterfactual
  wholesale cost on common-feasible days.
  \item The coherent-versus-independent feasible-day difference is statistically
  unresolved.
  \item The ex-post correction does not reverse the ranking.
  \item The result is bounded to SPARK, the frozen F1 asset, and
  \status{POST_PUBLICATION_TIMING_ASSUMED}.
\end{itemize}
\textbf{Forbidden claims}
\begin{itemize}
  \item Actual customer savings, trading profits, or plant validation.
  \item Honda economic value, general industrial-flexibility value, or battery
  value.
  \item Universal superiority or inferiority of coherent scenarios.
  \item That better dependence always improves decisions, or that it never does.
\end{itemize}}

\ifplanningappendix
\clearpage
\appendix

\section{Paper Outline and Target Allocation}
\label{app:outline}

Table~\ref{tab:outline} is an authoring blueprint for a 6--8 page applied
conference paper. It should be removed from the submitted manuscript by setting
\verb|\planningappendixfalse|.

\begin{table*}[h]
\centering
\caption{Section plan, target word counts, and exact message.}
\label{tab:outline}
\small
\begin{tabularx}{\textwidth}{p{0.23\textwidth}rX}
\toprule
Section & Target words & Exact message\\
\midrule
Abstract & 200--250 & Better temporal dependence may improve modelled feasibility without improving mean day-ahead decision cost.\\
1. Introduction & 400--500 & Decision value, rather than forecast score alone, is the scientific object; the central question is narrow and falsifiable.\\
2. Related work & 450--550 & Existing strands cover forecasting, calibration, scenarios, decision-focused methods, flexibility, and markets, but their integration requires direct decision evaluation.\\
3. Empirical foundation & 500--600 & SPARK demand and Honda signed net exchange provide distinct, leakage-controlled evidence; they must not be treated as equivalent targets.\\
4. Scenario and decision methods & 650--750 & S1 and S2 differ only in temporal dependence; all policies use one frozen counterfactual F1 asset and the same assumed information timing.\\
5. Evaluation protocol & 350--450 & The primary estimand is the 191-day common-feasible intersection; bootstrap, block-bootstrap, and correction rules are frozen and transparent.\\
6. Results & 700--850 & P4 reduces violations, P2 remains cheaper, and P3 versus P4 is unresolved; correction preserves the ranking.\\
7. Discussion & 450--600 & Dependence quality, feasibility, and mean economic value are distinct; mathematical structure may suppress the value of coherence.\\
8. Limitations and conclusion & 250--350 & The result is exploratory and bounded to one counterfactual asset, site, timing convention, and tariff setting.\\
\midrule
Total main text & 3,900--4,850 & Suitable for a concise 6--8 page two-column applied conference paper, subject to venue template and figure count.\\
\bottomrule
\end{tabularx}
\end{table*}

\section{Related-Work Writing Plan}
\label{app:related}

\begin{enumerate}
  \item \textbf{Industrial electricity forecasting:} establish day-ahead
  industrial load and net-exchange forecasting practice, data scarcity, and
  operational information cutoffs. Do not present Honda and SPARK as equivalent.
  \item \textbf{Probabilistic forecasting and quantile calibration:} distinguish
  marginal calibration, sharpness, and proper scoring from joint trajectory
  validity.
  \item \textbf{Scenario generation and temporal dependence:} cover residual
  bootstrap, trajectory sampling, energy score, variogram score, and dependence
  diagnostics. Explain why improved dependence does not guarantee decision value.
  \item \textbf{Decision-focused learning and predict-then-optimize:} position
  direct downstream evaluation without claiming to introduce a new
  decision-focused learning algorithm.
  \item \textbf{Industrial demand-side flexibility:} review process flexibility,
  virtual buffers, and physical constraints; label F1 as a hypothetical
  archetype rather than calibrated equipment.
  \item \textbf{Electricity-market timing and settlement:} verify gate closure,
  price publication, product resolution, and settlement conventions from
  authoritative sources.
  \item \textbf{Gap:} forecasting and scenario papers often stop at statistical
  scores, whereas this paper tests whether temporal coherence changes a frozen,
  feasibility-constrained downstream decision.
\end{enumerate}

\section{Novelty Positioning}
\label{app:novelty}

\begin{table*}[h]
\centering
\caption{Conservative novelty positioning pending independent literature review.}
\small
\begin{tabularx}{\textwidth}{p{0.23\textwidth}p{0.17\textwidth}XX}
\toprule
Candidate contribution & Classification & How the paper should claim it & What the paper must not claim\\
\midrule
Leakage-safe forecasting and provenance controls & Applied & A necessary empirical foundation for a timing-correct decision study & A new general forecasting methodology\\
Whole-day residual trajectory sampling & Established / applied & A reproducible implementation used to isolate temporal dependence & Invention of coherent scenario generation\\
Feasibility-restricted policy comparison & Useful applied contribution & A controlled comparison on identical feasible days with an explicit all-day correction & Universal causal identification of scenario value\\
Violation reduction without lower mean cost & Distinctive empirical finding & A bounded negative result under SPARK, F1, and the frozen timing/tariff setting & Proof that coherent scenarios are generally useful or useless\\
Positive-demand versus signed-exchange evidence & Applied cross-dataset context & Evidence that target semantics affect evaluation and interpretation & A common economic comparison or Honda decision-value result\\
Virtual production-buffer formulation & Established abstraction & A transparent hypothetical flexibility archetype & Plant-calibrated equipment, a battery, or reconstructed factory operations\\
Post-publication scheduling convention & Unsupported operationally & An explicitly assumed timing used for a counterfactual & Independently verified live-market implementability\\
Customer or trading value & Unsupported & Do not claim & Actual bill savings, revenue, profit, or deployable trading performance\\
\bottomrule
\end{tabularx}
\end{table*}

\section{Target-Venue Matrix}
\label{app:venues}

\begin{table*}[h]
\centering
\caption{Provisional venue positioning; current calls and formats must be verified.}
\small
\begin{tabularx}{\textwidth}{p{0.18\textwidth}p{0.22\textwidth}p{0.18\textwidth}p{0.18\textwidth}X}
\toprule
Venue & Fit & Audience & Likely acceptance risk & Required emphasis\\
\midrule
The smarter E Europe & Strong applied/industry fit if a scientific conference track is available & Energy managers, technology providers, flexibility and market practitioners & Academic depth or archival status may be insufficient; current format must be verified & Operational meaning, transparent limitations, feasibility, and actionable implications rather than model novelty\\
pro flex Dresden & Strong topical fit for industrial flexibility and grid interaction & German flexibility researchers, utilities, aggregators, and industry & Counterfactual asset and assumed timing may limit claims & Physical interpretation, constraint behavior, negative result, and German-market relevance\\
Energy Informatics & Strong methodological fit for data-to-decision evaluation & Energy-informatics and applied ML researchers & Requires verified literature novelty, reproducible methods, and possibly greater paper depth & Leakage safety, scenario diagnostics, decision-value separation, and open reproducibility\\
ETG/VDE energy event & Good power-systems and operations fit, depending on the specific call & Power engineers, system operators, utilities, and applied researchers & Hypothetical asset and absence of plant validation may be scrutinized & Units, market timing, feasibility, solver formulation, and restrained system-level conclusions\\
\bottomrule
\end{tabularx}
\end{table*}

\section{Reference Plan: No Citations Yet}
\label{app:references}

No bibliographic citation should enter the paper until its metadata, primary
source, relevance, and claim support have been checked. Every item below is
therefore marked \tverify{}.

\begin{table*}[h]
\centering
\caption{Reference acquisition and verification plan.}
\small
\begin{tabularx}{\textwidth}{p{0.22\textwidth}p{0.12\textwidth}X X}
\toprule
Reference type & Status & Evidence needed & Intended use\\
\midrule
Proper scoring rules and probabilistic forecast evaluation & \tverify & Foundational peer-reviewed source defining propriety and forecast-distribution evaluation & Separate marginal forecast quality from decision value\\
Weighted interval score and interval coverage & \tverify & Primary methodological source and accepted forecast-evaluation guidance & Define WIS and coverage without overstating joint calibration\\
Energy score & \tverify & Foundational source defining the multivariate proper score and known sensitivity limitations & Justify trajectory-level distribution assessment\\
Variogram score & \tverify & Foundational source defining the score and its sensitivity to dependence & Evaluate temporal dependence representation\\
Residual and block bootstrap & \tverify & Foundational statistical references for residual resampling and dependent-data bootstrap & Support scenario generation and paired/block uncertainty intervals\\
Time-series scenario generation & \tverify & Recent peer-reviewed energy-forecasting literature, preferably 2022--2026, plus foundational work & Position independent and coherent residual trajectories\\
Decision-focused learning / predict-then-optimize & \tverify & Foundational and recent survey or benchmark papers & Explain why forecast scores need not map to decision value\\
Value of probabilistic forecasts & \tverify & Peer-reviewed power/energy studies evaluating downstream operational value & Identify what is established and what this study adds empirically\\
Industrial demand-side flexibility & \tverify & Reviews and process-flexibility formulations with explicit physical constraints & Position F1 as an abstraction and bound physical claims\\
Virtual battery / production-buffer models & \tverify & Original or authoritative formulations distinguishing process buffers from electrochemical storage & Define state dynamics without calling F1 a real battery\\
Electricity-market design and day-ahead timing & \tverify & Official EPEX/SDAC, regulatory, or market-operator documents for the relevant year & Verify gate closure, price publication, product resolution, and DST semantics\\
Industrial tariffs and import/export remuneration & \tverify & Official tariff/regulatory documentation or explicit scenario literature & Prevent wholesale-price exposure from being called a customer bill\\
KIT SPARK dataset & \tverify & Official dataset paper, repository, DOI, license, and measurement definitions & Establish source identity and permitted claims\\
Honda dataset & \tverify & Official dataset paper, repository, DOI, license, and target/component definitions & Establish signed-grid-exchange semantics and provenance limits\\
Foundation time-series models & \tverify & Official model papers, repositories, checkpoint licenses, and versions & Provide bounded context only; not frame the paper around model ranking\\
\bottomrule
\end{tabularx}
\end{table*}

\section{Submission Package}
\label{app:submission}

\subsection{150-word conference abstract}

Industrial flexibility scheduling increasingly uses probabilistic forecasts,
but improved scenario dependence need not improve decisions. We compare
deterministic LightGBM scheduling with independent interval-residual and
coherent whole-day residual scenarios under a frozen hypothetical
electrical-work buffer. The leakage-safe counterfactual evaluation uses 273
complete 15-minute KIT SPARK days in 2024; 191 days are feasible for all
policies. On this common-feasible population, deterministic scheduling is
cheaper than coherent scheduling by 0.191191 EUR/day, while the
independent-versus-coherent difference is statistically unresolved. Coherent
scenarios nevertheless reduce modelled zero-export-cap violations from 82
deterministic days and 16 independent-scenario days to 11. An explicit all-day
feasibility correction does not reverse the cost ranking. The findings show
that improved temporal dependence can enhance modelled feasibility without
reducing mean synthetic counterfactual wholesale cost. Results are exploratory,
rely on a hypothetical asset and assumed post-publication timing, and
constitute neither plant validation, customer savings, nor trading profit.
This uses the frozen risk-neutral expected-cost model.

\subsection{250-word extended abstract}

Probabilistic forecasts are often evaluated using calibration and multivariate
scores, although a better representation of temporal dependence does not
necessarily improve the downstream schedule. We test whether coherent
whole-day residual trajectories add decision value relative to a deterministic
LightGBM forecast and independent interval-residual scenarios. All policies use
a hypothetical electrical-work buffer with $\pm5$ kW power,
$\pm10$ kWh state capacity, terminal recovery, and a modelled zero-export
cap. Forecasts are issued at 12:00 Europe/Berlin on D-1; price-aware scheduling
at 12:55 is assumed rather than verified from observation-level publication
timestamps.

The counterfactual evaluation contains 273 complete 15-minute KIT SPARK days
from January to December 2024. The primary paired estimand uses 191 days that
are feasible for deterministic, independent-scenario, and coherent-scenario
policies. Deterministic scheduling is cheaper than coherent scheduling by
0.191191 EUR/day (95\% whole-day bootstrap interval: -0.271593 to -0.131960
for deterministic minus coherent cost). It is also cheaper than independent-
scenario scheduling by 0.143951 EUR/day. Independent versus coherent scheduling
is unresolved: the mean difference is -0.047240 EUR/day and its 95\% interval
crosses zero. Coherent scheduling nevertheless reduces modelled zero-export-cap
violations from 82 deterministic days and 16 independent-scenario days to 11.
An external feasibility correction over all 273 days does not reverse
the cost ranking.

Thus, improved temporal dependence changes feasibility without reducing mean
synthetic counterfactual wholesale cost. The result is exploratory and bounded
to one dataset, a hypothetical asset, wholesale exposure, imperfect
marginal calibration, and assumed post-publication timing. It does not establish
plant performance, customer savings, trading profit, or universal inferiority
of coherent scenarios.

\subsection{Presentation learning outcomes}

After the presentation, participants should be able to:
\begin{enumerate}
  \item distinguish marginal calibration, temporal coherence, modelled
  feasibility, and realized decision value;
  \item explain why a risk-neutral expected-cost problem may not reward more
  realistic residual dependence; and
  \item design a common-feasible, pre-registered comparison that reports cost
  and constraint violations without overstating counterfactual outcomes.
\end{enumerate}

\subsection{Likely reviewer questions and concise answers}

\begin{enumerate}
  \item \textbf{Is the F1 resource a real industrial asset?} No. It is a frozen
  hypothetical electrical-work buffer used to isolate decision effects; its
  capacities are not plant measurements.
  \item \textbf{Why exclude infeasible days from the primary comparison?} The
  191-day intersection ensures identical support and physically comparable
  schedules. A pre-specified corrected all-day analysis is reported separately.
  \item \textbf{Does P4's lower violation rate prove economic value?} No. It is
  a feasibility result. P4 has higher mean cost than P2, and P3 versus P4 is
  unresolved on common-feasible days.
  \item \textbf{Is the 12:55 decision time operationally verified?} No. It is
  explicitly labelled \status{POST_PUBLICATION_TIMING_ASSUMED} because the
  historical price file lacks observation-level publication timestamps.
  \item \textbf{Can the result be generalized to Honda or batteries?} No.
  Honda has signed net-exchange semantics and no economic evaluation here; F1
  is neither a battery nor validated production equipment.
\end{enumerate}

\subsection{One-slide paper narrative}

\begin{description}
  \item[Problem.] Industrial schedules increasingly consume probabilistic
  forecasts, but improved scenario diagnostics may not improve decisions.
  \item[Gap.] Temporal coherence is rarely isolated from physical feasibility
  and evaluated on identical downstream decision days.
  \item[Method.] Compare deterministic P2, independent-scenario P3, and coherent-
  scenario P4 schedules with one frozen hypothetical F1 asset.
  \item[Data.] 273 complete SPARK days in 2024; 191 are feasible for all three
  policies. Honda supplies bounded cross-dataset forecasting and scenario
  evidence, not economics.
  \item[Result.] P4 reduces modelled violation days to 11, but P2 is cheaper by
  0.191191 EUR/day on common-feasible days; P3 versus P4 is unresolved.
  \item[Implication.] Better temporal dependence is not equivalent to better
  decision value; alignment with the economic structure must be tested.
\end{description}

\subsection{Suggested 10-minute presentation outline}

\begin{enumerate}
  \item \textbf{0:00--0:45 --- Question and claim boundary.} State the decision
  question and that all outcomes are counterfactual.
  \item \textbf{0:45--1:45 --- Data and leakage controls.} Explain SPARK demand,
  Honda signed exchange, provenance, and timing.
  \item \textbf{1:45--3:00 --- Scenario methods.} Contrast independent S1 and
  whole-day S2 using dependence and coverage diagnostics.
  \item \textbf{3:00--4:15 --- Frozen F1 decision problem.} Present state,
  power, terminal, and zero-export constraints and the assumed timeline.
  \item \textbf{4:15--6:15 --- Primary results.} Show common-feasible paired
  costs and bootstrap intervals.
  \item \textbf{6:15--7:15 --- Feasibility result.} Show violation frequency,
  overlap, and external correction.
  \item \textbf{7:15--8:30 --- Mechanism.} Explain marginal expectations,
  conservatism, and frequent binding constraints.
  \item \textbf{8:30--9:30 --- Limitations and practice.} Bound asset, timing,
  tariff, exposed test year, and generalizability.
  \item \textbf{9:30--10:00 --- Takeaway.} Scenario sophistication should match
  the downstream economic structure.
\end{enumerate}

\subsection{Suggested slide-deck title}

\begin{center}
\emph{When Scenario Coherence Improves Feasibility but Not Cost: Evidence from
Counterfactual Industrial Flexibility Scheduling}
\end{center}

\subsection{Author biography template}

\textbf{[Author Name]} is an MSc Digital Engineering student at
Bauhaus-Universit\"at Weimar. [He/She/They] previously worked in energy-utility
analytics, with experience in [verified domain or role]. [His/Her/Their]
research interests include industrial energy forecasting, uncertainty-aware
decision support, and reproducible optimization under operational constraints.
This biography should be completed only with verifiable education and
professional details.

\section{Author Checklist Before Submission}

\begin{itemize}
  \item Replace every \tverify{} literature item with a checked primary source;
  do not invent or infer citation metadata.
  \item Convert the source to the selected venue's official template and confirm
  its current page, anonymization, and reference requirements.
  \item Retain the terms ``counterfactual,'' ``modelled violation,'' and
  ``synthetic counterfactual wholesale-cost difference.''
  \item Keep all Honda evidence contextual; do not add Honda economic value.
  \item State the previously exposed 2024 evaluation period and the assumed
  market-publication timing.
  \item Add only figures whose source data and deterministic generation scripts
  are archived and traceable.
  \item Set \verb|\planningappendixfalse| before producing the submission PDF.
\end{itemize}
\fi

% Bibliography intentionally omitted. Add a verified .bib file only after the
% reference plan in Appendix E has been completed.

\end{document}
